\chapter{Preparing Your Disquisition Source Files}
\label{chap:preparing-source}

By the time you are preparing a disquisition, much of the difficult intellectual work is already behind you. This chapter focuses on the document task that remains: placing your content into the template in a clear, structured way. The goal is not to make you a \LaTeX{} expert. The goal is to help you work comfortably enough with the source files that the template can apply consistent formatting and produce a PDF that is easier to review.

This chapter explains the basic structure of the source files and introduces the most common kinds of content you are likely to add. The examples are intentionally brief. They are meant to show patterns, not every possible variation. Later chapters explain how to review the generated PDF and complete accessibility review and remediation workflows.

\section{Document information, front matter, and chapters}
\label{sec:document-information}

Most students begin by replacing the sample document information in \texttt{main.tex}. Depending on the repository version and the requirements of your program, this may include the disquisition title, author name, degree, program or department, committee information, date or term, and PDF metadata. The template controls the placement and formatting of this information. Your responsibility is to supply accurate content and remove placeholder text before final review.

Front matter files are stored in \texttt{frontmatter/}. These files may include the abstract, acknowledgments, dedication, abbreviations, or other introductory material. Some front matter sections are optional and are controlled by switches in \texttt{main.tex}. If a section is not needed, disable the corresponding switch rather than manually deleting the template structure.

The files in \texttt{chapters/} contain the main body chapters of the disquisition. The files in \texttt{appendices/} contain appendix material. Most students will spend the majority of their time working in these chapter and appendix files rather than editing the template itself. A chapter normally begins with real heading commands, such as \texttt{\textbackslash chapter\{Methods\}}, \texttt{\textbackslash section\{Data Collection\}}, and \texttt{\textbackslash subsection\{Interview Protocol\}}.
Use headings in a logical order. Do not create headings by manually changing font size, centering text, or typing numbers by hand. Real heading commands support the table of contents, bookmarks, cross-references, and PDF structure generated by the template.

Compile the main file, not an individual chapter file. Chapter files depend on the main file for document setup, template behavior, bibliography settings, and overall structure.

\section{Figures, tables, equations, and references}
\label{sec:common-content}

Figures, tables, equations, citations, and links are easiest to manage when they are treated as structured content rather than manually formatted objects. The template includes commands and environments for these common patterns.

Store images in the repository's \texttt{figures/} folder or in the image location expected by the project. Meaningful figures should have captions and descriptions. The caption identifies the figure for all readers. The description or alternative text explains the purpose or important information contained in the image. A typical figure pattern in this template looks like this:
\begin{ndsucodelisting}[Figure code example]
\begin{figure}[htbp]
\centering

\ndsuincludegraphics[
    width=\textwidth,
    keepaspectratio
]
{Alternative text describing the image.}
{figures/example-image.png}

\caption{Example figure caption.}
\label{fig:example}

\ndsufigurenote{Optional figure note.}

\end{figure}
\end{ndsucodelisting}

The exact image, caption, and description should match the content of your disquisition. Decorative images should use the decorative-image workflow provided by the template rather than being treated as meaningful figures.

Tables should present structured information rather than visual layout. Keep tables as simple as the content allows, use clear headers, and avoid screenshots of tables when possible. A screenshot is an image, not structured table data. Very wide or complicated tables may need to be simplified, split into smaller parts, or moved to an appendix.

Equations should use \LaTeX{} math rather than images of equations. Use inline math when notation belongs inside a sentence (e.g., a short equation like $e = mc^2$). Inline notation is useful for brief definitions, but longer or more important equations should usually be displayed and numbered when you need to refer to them later. Equation~\ref{eq:state-cost-example} expands the inline state update into a compact display equation. It is included to demonstrate equation numbering, matrices, fractions, Greek letters, superscripts, subscripts, an integral, and a List of Equations entry when that list is enabled. In your own disquisition, use nearby prose to explain what the variables mean and why the equation matters for your argument. Complex mathematical content may need additional review after the PDF is generated.

\equationlistentry{Example state update and cost function}
\mathalt{The next state s sub t plus one equals a two by two transition matrix times the current state s sub t plus a two element input vector times u sub t. The cost J of s equals the integral from zero to T of s transpose Q s plus rho times u squared, with respect to t.}
\begin{equation}
\begin{aligned}
  s_{t+1}
    &=
    \begin{bmatrix}
      1 & \Delta t \\
      0 & 1
    \end{bmatrix}
    s_t
    +
    \begin{bmatrix}
      \frac{1}{2}\Delta t^2 \\
      \Delta t
    \end{bmatrix}
    u_t,
  \\[0.5em]
  J(s)
    &=
    \int_{0}^{T}
    \left(
      s(t)^{\mathsf{T}} Q s(t)
      + \rho u(t)^2
    \right)
    \, dt .
\end{aligned}
\label{eq:state-cost-example}
\end{equation}


Citations and references depend on two parts working together: a bibliography entry in \texttt{references.bib} and a citation command in the chapter text. The template uses the \texttt{biblatex} workflow with Biber. If a citation appears as question marks or does not update, check the citation key, the bibliography entry, and the compile log before changing unrelated source files. Use descriptive link text rather than bare URLs or repeated phrases such as ``click here.'' A reader should be able to understand where a link goes from the linked text itself.

\section{Working with the source and reviewing changes}
\label{sec:working-reviewing}

A useful workflow is to make one meaningful change at a time, compile, and review the result. This is especially helpful when adding figures, tables, equations, citations, or appendices. If something fails, the source file, the compile log, and the most recent change usually provide the fastest path to diagnosis.
When the PDF does not look right, fix the source first whenever possible. Do not treat the generated PDF as the place to rebuild headings, renumber chapters, manually correct references, or redesign front matter. Source-level fixes are easier to repeat and are less likely to be lost when the PDF is regenerated. Before compiling a more complete review build, use Table~\ref{tab:source-review-checks} to check the source for common unfinished items. This table is not a final submission checklist. It is a practical source check before you move into PDF review.

\begin{table}[htbp]
\centering
\caption{Source checks before compiling a review PDF.}
\label{tab:source-review-checks}

\begin{ndsudatatable}[headers={1}]
\begin{tabularx}{\textwidth}{L{1.55in}XX}
\toprule
\textbf{Source check} & \textbf{Example} & \textbf{What to fix before review} \\
\midrule
Document information & Title, author, degree, program, committee names, dates, subject, and keywords & Replace sample names, sample titles, placeholder dates, and unfinished metadata in \texttt{main.tex}. \\

Chapters and appendices & Chapter files in \texttt{chapters/} and appendix files in \texttt{appendices/} & Confirm that each intended file is included once and that unused sample files are removed or disabled. \\

Headings & \texttt{\textbackslash chapter\{\}}, \texttt{\textbackslash section\{\}}, and \texttt{\textbackslash subsection\{\}} headings & Use real heading commands in a logical order instead of bold text, manual numbering, or extra spacing. \\

Figures & Meaningful images, diagrams, screenshots, charts, and maps & Confirm that figures have captions, labels when referenced, readable image quality, and source-provided descriptions. \\

Tables & Data tables and review tables & Use real table structure with clear headers. Avoid screenshots of tables when structured table data would work better. \\

Equations & Numbered equations and technical notation & Introduce important equations in prose, define variables when needed, and label equations that are referenced later. \\

Citations and references & In-text citations connected to \texttt{references.bib} & Check that citation keys match bibliography entries and that the References section is not empty unless that is intentional. \\

Links & External links and internal cross-references & Use meaningful link text and confirm that cross-references point to the intended figure, table, equation, chapter, or appendix. \\

Wide or landscape content & Wide tables, large figures, diagrams, schemas, or rotated pages & Decide whether the content should be simplified, split, moved to an appendix, or flagged for closer PDF review. \\
\bottomrule
\end{tabularx}
\end{ndsudatatable}

\ndsutablenote{Use this table as a source-review aid before compiling a more complete PDF. Final formatting, accessibility review, and any PDF remediation still require review of the generated PDF.}
\end{table}

The next chapter explains how to compile and review the generated PDF, including checks for figures, tables, equations, citations, links, and landscape pages.
